\begin{initiativkarte}{Computer Science for Future (CS4F)}

  % Bild (optional)
  \IfFileExists{bilder/cs4f.jpg}{%
    \begin{center}
      \includegraphics[width=\linewidth,height=5cm,keepaspectratio]{bilder/cs4f.jpg}
    \end{center}
    \vspace{0.4cm}
  }{}

  % Kurzbeschreibung
  \textit{\textcolor{hawgray}{Computer Science for Future (CS4F) beschäftigt sich mit nachhaltiger Informatik und der Integration von Klimaschutz, Ethik und gesellschaftlicher Verantwortung in Lehre, Forschung und Organisation.}}

  \vspace{0.5cm}
  \hawrule

  % Beschreibung
  \textbf{\textcolor{hawblue}{Was machen wir?}}\\[0.3cm]

  Computer Science for Future (CS4F) ist eine Initiative des Departments Informatik der HAW Hamburg, die sich mit nachhaltiger Entwicklung in der Informatik beschäftigt.

  Grundlage ist die Agenda 2030 der Vereinten Nationen mit ihren Sustainable Development Goals (SDGs), die ökologische, ökonomische und soziale Nachhaltigkeit miteinander verbinden.

  CS4F verfolgt das Ziel, Informatik nicht nur effizienter, sondern gesellschaftlich verantwortlicher zu gestalten. Dabei geht es darum, von einem reinen „Schaden reduzieren“ hin zu echter aktiver Verantwortung in der Gestaltung digitaler Systeme zu gelangen.

  Die Initiative wirkt auf mehreren Ebenen:
  \begin{itemize}
      \item Integration von Nachhaltigkeit in die Lehre
      \item Zusammenarbeit mit Forschung und Industrie
      \item Reflexion ethischer Fragestellungen in der Informatik
  \end{itemize}

  Aktuelle Lehrangebote und Projekte behandeln unter anderem:
  \begin{itemize}
      \item Menschenrechte und digitale Protokolle
      \item Klimawandel und IT-Systeme
      \item Maschinenethik und autonome Systeme

  \end{itemize}


  % Tags
  \tag{Nachhaltigkeit}
  \tag{Informatik}
  \tag{Ethik}
  \tag{Forschung}
  \tag{Interdisziplinär}

  \vspace{0.6cm}
  \hawrule

  % Infozeilen
  \infozeile{\faUsers}{Zielgruppe: Studierende der Informatik und angrenzender Fachrichtungen}
  \infozeile{\faGraduationCap}{Semester: Ab dem 1. Semester (Vertiefung später möglich)}
  \infozeile{\faMapMarker*}{Ort: HAW Hamburg Department Informatik}
  \infozeile{\faGlobe}{Website: \href{https://www.haw-hamburg.de/hochschule/technik-und-informatik/departments/informatik/unser-department/computer-science-for-future/}{CS4F Projektseite}}

\end{initiativkarte}